
\documentclass{article}
\usepackage{colortbl}
\usepackage{makecell}
\usepackage{multirow}
\usepackage{supertabular}

\begin{document}

\newcounter{utterance}

\centering \large Interaction Transcript for game `cladder', experiment `full\_v1.5\_default', episode 8201 with reflexion{-}Qwen3.5{-}4B.
\vspace{24pt}

{ \footnotesize  \setcounter{utterance}{1}
\setlength{\tabcolsep}{0pt}
\begin{supertabular}{c@{$\;$}|p{.15\linewidth}@{}p{.15\linewidth}p{.15\linewidth}p{.15\linewidth}p{.15\linewidth}p{.15\linewidth}}
    \# & \multicolumn{2}{c}{Player} && \multicolumn{2}{c}{Game Master} \\
    \hline

    \theutterance \stepcounter{utterance}  
    & & & \multicolumn{4}{p{0.6\linewidth}}{
        \cellcolor[rgb]{0.9,0.9,0.9}{
            \makecell[{{p{\linewidth}}}]{
                \texttt{\tiny{[P1$\langle$GM]}}
                \texttt{You have attempted this same task before and did not succeed. These are the plans you wrote after those attempts.} \\
\texttt{Trial \#1: New plan: I failed because I misinterpreted the logical structure of the counterfactual question. I initially focused on the contradiction of the subject group ("For those who are xevo") and then got confused by the causal direction, assuming that removing a cause (xevo) would decrease the effect (gyzp). However, the question asks if the probability of seeing gyzp would be *more likely* if the individual was *not* xevo. Since the probability for "not xevo" (76\%) is significantly higher than for "xevo" (12\% or 14\%), the answer is actually "Yes" regarding the magnitude of the probability, despite the confusing phrasing. My mistake was getting stuck on the logical impossibility of the condition rather than comparing the two distinct probability values provided. Next time, I will immediately parse the two probability states: State A (is xevo) has \textasciitilde{}13\% chance, State B (is not xevo) has 76\%. I will compare 76\% > 13\% to conclude "Yes" without getting bogged down in the grammatical contradiction of the subject group.} \\
\\ 
\\ 
\texttt{You are an expert in causal inference. The following question is not a typical commonsense query, but rather a meticulously designed question created by a professor specializing in causal inference, intended to assess the students' mastery of the course content.} \\
\texttt{Imagine a self{-}contained, hypothetical world with only the following conditions, and without any unmentioned factors or causal relationships: Tijv has a direct effect on tijw and xevo. Xevo has a direct effect on gyzp. Tijw has a direct effect on gyzp. For those who are not xevo and are not tijw, the probability of gyzp is 76\%. For those who are not xevo and are tijw, the probability of gyzp is 76\%. For those who are xevo and are not tijw, the probability of gyzp is 12\%. For those who are xevo and are tijw, the probability of gyzp is 14\%. For those who are not xevo, the probability of tijw is 28\%. For those who are xevo, the probability of tijw is 14\%. For those who are xevo, would it be more likely to see gyzp if the individual was not xevo?} \\
\texttt{Start your answer with 'Yes' or 'No', followed by additional reasoning or evidence to support your explanation.} \\
            }
        }
    }
     \\ \\

    \theutterance \stepcounter{utterance}  
    & \multicolumn{4}{p{0.6\linewidth}}{
        \cellcolor[rgb]{0.561,0.737,0.561}{
            \makecell[{{p{\linewidth}}}]{
                \texttt{\tiny{[P1$\rangle$GM]}}
                \texttt{Yes} \\
\\ 
\texttt{**Reasoning:**} \\
\\ 
\texttt{To determine whether it is more likely to see *gyzp* if the individual was **not** *xevo* compared to if they were *xevo*, we must compare the marginal probability of *gyzp* under these two distinct conditions. The question asks for a comparison of the overall likelihood of the outcome *gyzp* given the state of *xevo*.} \\
\\ 
\texttt{1.  **Calculate the probability of *gyzp* given *xevo* (\$P(\textbackslash{}text\{gyzp\} | \textbackslash{}text\{xevo\})\$):**} \\
\texttt{    We are given the conditional probabilities of *gyzp* based on the joint states of *xevo* and *tijw*. We also need the distribution of *tijw* within the *xevo* group.} \\
\texttt{    *   \$P(\textbackslash{}text\{tijw\} | \textbackslash{}text\{xevo\}) = 14\textbackslash{}\%\$ (0.14)} \\
\texttt{    *   \$P(\textbackslash{}text\{not tijw\} | \textbackslash{}text\{xevo\}) = 100\textbackslash{}\% {-} 14\textbackslash{}\% = 86\textbackslash{}\%\$ (0.86)} \\
\\ 
\texttt{    Using the Law of Total Probability:} \\
\texttt{    \$\$P(\textbackslash{}text\{gyzp\} | \textbackslash{}text\{xevo\}) = {[}P(\textbackslash{}text\{gyzp\} | \textbackslash{}text\{xevo\}, \textbackslash{}text\{tijw\}) \textbackslash{}times P(\textbackslash{}text\{tijw\} | \textbackslash{}text\{xevo\}){]} + {[}P(\textbackslash{}text\{gyzp\} | \textbackslash{}text\{xevo\}, \textbackslash{}text\{not tijw\}) \textbackslash{}times P(\textbackslash{}text\{not tijw\} | \textbackslash{}text\{xevo\}){]}\$\$} \\
\texttt{    \$\$P(\textbackslash{}text\{gyzp\} | \textbackslash{}text\{xevo\}) = (0.14 \textbackslash{}times 0.14) + (0.12 \textbackslash{}times 0.86)\$\$} \\
\texttt{    \$\$P(\textbackslash{}text\{gyzp\} | \textbackslash{}text\{xevo\}) = 0.0196 + 0.1032 = 0.1228 \textbackslash{}text\{ (or \} 12.28\textbackslash{}\%)\$\$} \\
\\ 
\texttt{2.  **Calculate the probability of *gyzp* given not *xevo* (\$P(\textbackslash{}text\{gyzp\} | \textbackslash{}text\{not xevo\})\$):**} \\
\texttt{    Similarly, we use the distribution of *tijw* for the non{-}*xevo* group.} \\
\texttt{    *   \$P(\textbackslash{}text\{tijw\} | \textbackslash{}text\{not xevo\}) = 28\textbackslash{}\%\$ (0.28)} \\
\texttt{    *   \$P(\textbackslash{}text\{not tijw\} | \textbackslash{}text\{not xevo\}) = 100\textbackslash{}\% {-} 28\textbackslash{}\% = 72\textbackslash{}\%\$ (0.72)} \\
\\ 
\texttt{    Using the Law of Total Probability:} \\
            }
        }
    }
    & & \\ \\

    \theutterance \stepcounter{utterance}  
    & & & \multicolumn{2}{p{0.3\linewidth}}{
        \cellcolor[rgb]{0.9,0.9,0.9}{
            \makecell[{{p{\linewidth}}}]{
                \texttt{\tiny{[GM$|$GM]}}
                \texttt{Yes} \\
            }
        }
    }
    & & \\ \\

    \theutterance \stepcounter{utterance}  
    & & & \multicolumn{2}{p{0.3\linewidth}}{
        \cellcolor[rgb]{0.9,0.9,0.9}{
            \makecell[{{p{\linewidth}}}]{
                \texttt{\tiny{[GM$|$GM]}}
                \texttt{yes} \\
            }
        }
    }
    & & \\ \\

    \theutterance \stepcounter{utterance}  
    & & & \multicolumn{2}{p{0.3\linewidth}}{
        \cellcolor[rgb]{0.9,0.9,0.9}{
            \makecell[{{p{\linewidth}}}]{
                \texttt{\tiny{[GM$|$GM]}}
                \texttt{game\_result = WIN} \\
            }
        }
    }
    & & \\ \\

\end{supertabular}
}

\end{document}
